\section{Introduction}\label{sec:v7-introduction}
A statistic that preserves present predictions need not admit a finite
recursive implementation. Conversely, a small implementable register
need not preserve the decision problems of the underlying experiment.
The purpose of this paper is to relate these three objects: a predictive
quotient, its attainable geometry, and a charged causal realization.
The order of the argument is
\[
 \begin{gathered}
 \text{positive causal experiment}\longrightarrow\text{future-test quotient}\\
 \longrightarrow\text{causal simulation}\longrightarrow\text{finite-resource realization}.
 \end{gathered}
\]
The quotient is always relative to a declared family of executable future
tests. Resource preservation is always relative to a declared resource
signature. Neither a measurable re-encoding by real numbers nor a
terminal randomization is, by itself, a finite-memory implementation.

The principal structural result is Theorem~\ref{thm:v7-tree}. It starts
with a factorial language, a stationary law, and a positive energy on its
finite words. The hypotheses are local child comparisons, strict energy
decrease under prefix attachment, bounded unary chains, and geometric
separation of the predictive images. No quantization dimension is an
input. If $G_M$ is the energy of the greedy cut with at most $M$ leaves,
then its static squared quantization error is comparable to $G_M$.
Every greedy tree is suffix closed. With $L$ leaves it has at most
$(U+1)(2L-1)$ vertices, where $U$ bounds unary chains. These vertices
support two different causal operations. Deleting the oldest available
future symbol realizes a renewal orbit code; prepending the newest
observed symbol realizes a code of the entire observed past. Both
operations have the same-order risk as the unrestricted static problem.
The lower bounds allow randomized, time-dependent observers.

The two orientations are not merely two descriptions of one reset model.
The Markov-renewal verification in Section~\ref{sec:v7-instances} includes
forbidden transitions and noninjective, delay-identifying observations.
The second verification is a noisy hidden shift register whose tail
persists at every time. Its complete posterior has infinitely many
observation-dependent coordinates. Theorem~\ref{thm:v7-shift} gives a
matched whole-register excess-risk law in that experiment, including its
noise dependence. Thus one structural realization theorem has a
regenerative and a genuinely continually observed application.

A complementary approximation theorem treats predictive quotients that
are measure valued rather than tree separated.
Theorem~\ref{thm:v7-filter} constructs a finite recursive register for
contracting filters on $\mathcal P([0,1])$. Its codebook has at most
$4^n$ states and error at most $3/(2n(1-\rho))$, apart from initialization
and explicitly charged numerical defects. The nonlinear example has no
fixed-step common minorizer and no closed recursion of any prescribed
finite list of polynomial moments. This is an upper approximation
theorem, not a claim that every nonlinear filter has this contraction or
that its logarithmic budget rate is universally optimal.

The experiment-comparison application is two-sided.
Theorem~\ref{thm:v7-gaussian-budget} proves that an $r$-dimensional
nondegenerate Gaussian location experiment, over a compact parameter
set containing a ball, has optimal $M$-outcome approximation error
$\Theta(M^{-2/r})$ in Le Cam distance. The upper bound uses deterministic
cells and positive second-order reconstruction, not first-order constant
histograms. The lower bound applies to every experiment on at most $M$
outcomes. For the actual local count-to-contact protocol of A2,
Theorem~\ref{thm:v7-a2} proves
\[
 \Del(\mathcal E_{N,J},\mathcal G)
 \le C\{N^{-1/2}+J N^{-1/2}+J^{-2}\},\qquad
 |\mathcal E_{N,J}|\le J^r.
\]
It follows that $\Del\asymp M^{-2/r}$ throughout the growing-budget
window $1\ll M\lesssim N^{r/6}$. The proof includes the local normal,
variance replacement, boundary crossing, reconstruction, and decision
comparison arguments. The physical statistic is deterministic throughout.

\subsection{Relation to previous work}
The existence of path laws, conditional-probability quotients, Bayes
orthogonality, information loss under coarsening, and the composition of
statistical randomizations are classical; see
\cite{Kallenberg2021,Blackwell1953,LeCamYang2000,Torgersen1991}.
Sections~\ref{sec:v7-experiments}--\ref{sec:v7-transport} give the forms
needed here, with the simulator and its persistent state included. We do
not claim novelty for these identities or for their categorical notation.

Static graph-directed quantization is a particularly close antecedent.
Kesseb\"ohmer--Zhu \cite{KZ2017}, Theorem 1.1, determine Markov-type
quantization dimension by the spectral radius of a matrix of
$(p_{ij}c_{ij}^{r})^{s/(s+r)}$ and characterize finiteness of the upper
coefficient through maximal strongly connected components. Their
spatial theorem is not a theorem about a causal register. Our tree
argument separates that familiar static mechanism from the additional
suffix-closure condition and proves realization in both time
orientations. The Markov application retains the renewal-weighted orbit
and survival contribution; its spatial pressure is not a new discovery
of graph-directed quantization. The earlier comparisons with functional
quantization, dynamical rate distortion and zero-delay coding remain
relevant \cite{LuschgyPages2004,LinderYuksel2014,WoodLinderYuksel2017}.
A finite channel alphabet does not generally bound the entire persistent
state of its coding policy.

Demirci--Kara--Y\"uksel \cite{DKYAC2024}, Theorems 3.5--3.6 of the cited
preprint, obtain quantized-belief and finite-window policy approximations
in their contraction-based control analysis. Our measure-valued result
uses the classical contraction-plus-quantization mechanism, but provides
a concrete entire-register cardinality and a same-observation recursive
error bound for a specified nonlinear experiment. Its objective is
prediction and terminal experiment simulation, not average-cost control
optimality. The sharp hidden-shift theorem is instead a two-sided
predictive-tree result. Neither theorem purports to subsume general
partially observed control.

The new contribution is therefore not ``finite memory'', ``pressure'',
or ``Bayes reduction'' in isolation. It is the combination of a
checkable energy criterion, exact two-orientation realization, unrestricted
register converses, and resource-typed experiment comparison. The
Gaussian/A2 application supplies a separate decision-theoretic
cardinality obstruction and an attainable joint sample/budget law.
No priority claim for every component of these arguments is required.

\subsection{Reading convention}
This article is the canonical seventh revision of Volume I. Its proofs
are self-contained at the level of the stated models. A separate
\emph{Complete development} contains this article followed by the
unchanged earlier theorem bodies, including pressure, critical renewal,
calibration, collision, and control results. The companion is a
preservation and extended-reference object, not an alternative definition
of the main paper. Source identities, historical consultation, and
responses to the referee are recorded outside the mathematical text.
